% GENERATED FILE. Edit canonical lessons, then run npm run generate:books.
% canonical-source-hash: fnv1a64:ba265ecc95b47baa
% canonical-lessons: MW-C208-aspatal, MW-C208-tahsil, MW-C208-karkhano, MW-C208-adalat, MW-C208-khulno

\chapter{Offices in Town}
\label{ch:mw-offices-in-town}
\hlchaptermodality{208}

\begin{chapteropening}
\textbf{By the end of this chapter:} \emph{I can find an office or a hospital in town.}

Complete the last lesson of chapter 208: I can find an office or a hospital in town.
\end{chapteropening}

\section[\texorpdfstring{\mw{अस्पताल} (aspatāl)}{अस्पताल (aspatāl)}]{\mw{अस्पताल} (aspatāl) — a hospital}

\label{lesson:MW-C208-aspatal}



\begin{quote}
Before the new one: say the Marwadi for a big gate, then the Marwadi for a tower.
\end{quote}

\subsection*{You'll want to know: \mw{अस्पताल}}
\textbf{\mw{अस्पताल}} — \emph{aspatāl} — \textquotedblleft{}a hospital\textquotedblright{}.

\begin{quote}
\mw{अस्पताल} \mw{नदी} \mw{कनै} \mw{है।}
\end{quote}

\begin{quote}
\emph{aspatāl nadī kanai hai.}
\end{quote}

\begin{quote}
The hospital is near the river.
\end{quote}

A masculine noun.

\begin{cousinweb}
\textbf{\mw{अस्पताल}} (\emph{aspatāl}) is borrowed from English \emph{hospital}, reshaped in everyday speech.
\end{cousinweb}

\begin{grammarlens}[title={What you've built}]
One more word for this chapter.
\end{grammarlens}

\subsection*{Your turn}
\begin{itemize}
\raggedright
  \item \emph{Say it:} \emph{aspatāl}
  \item \emph{Say it:} \emph{aspatāl}, once more
  \item \emph{Say it:} \emph{aspatāl nadī kanai hai.} — The hospital is near the river.
  \item \emph{Recall:} say the Marwadi for a big gate, then the Marwadi for a tower, then say \emph{aspatāl} again
\end{itemize}

\begin{tcolorbox}[breakable,colback=teal!4,colframe=teal!35!black,title={Before you move on}]
What does \mw{अस्पताल} mean? (\textquotedblleft{}A hospital\textquotedblright{}.) Say it once more.
\end{tcolorbox}
\section[\texorpdfstring{\mw{तहसील} (tahsīl)}{तहसील (tahsīl)}]{\mw{तहसील} (tahsīl) — a sub-district office; a tehsil}

\label{lesson:MW-C208-tahsil}



\begin{quote}
Before the new one: say the Marwadi for a tower, then the Marwadi for a hospital.
\end{quote}

\subsection*{You'll want to know: \mw{तहसील}}
\textbf{\mw{तहसील}} — \emph{tahsīl} — \textquotedblleft{}a sub-district office; a tehsil\textquotedblright{}.

\begin{quote}
\mw{तहसील} \mw{सोमवार} \mw{बंद} \mw{है।}
\end{quote}

\begin{quote}
\emph{tahsīl somvār band hai.}
\end{quote}

\begin{quote}
The tehsil office is closed on Monday.
\end{quote}

A feminine noun.

\begin{cousinweb}
\textbf{\mw{तहसील}} (\emph{tahsīl}) is usually traced, through Persian, to Arabic \emph{taḥṣīl}, \textquotedblleft{}collecting\textquotedblright{}: probably the office where land revenue was collected.
\end{cousinweb}

\begin{grammarlens}[title={What you've built}]
One more word for this chapter.
\end{grammarlens}

\subsection*{Your turn}
\begin{itemize}
\raggedright
  \item \emph{Say it:} \emph{tahsīl}
  \item \emph{Say it:} \emph{tahsīl}, once more
  \item \emph{Say it:} \emph{tahsīl somvār band hai.} — The tehsil office is closed on Monday.
  \item \emph{Recall:} say the Marwadi for a tower, then the Marwadi for a hospital, then say \emph{tahsīl} again
\end{itemize}

\begin{tcolorbox}[breakable,colback=teal!4,colframe=teal!35!black,title={Before you move on}]
What does \mw{तहसील} mean? (\textquotedblleft{}A sub-district office; a tehsil\textquotedblright{}.) Say it once more.
\end{tcolorbox}
\section[\texorpdfstring{\mw{कारखानो} (kārkhāno)}{कारखानो (kārkhāno)}]{\mw{कारखानो} (kārkhāno) — a factory}

\label{lesson:MW-C208-karkhano}



\begin{quote}
Before the new one: say the Marwadi for a hospital, then the Marwadi for a sub-district office.
\end{quote}

\subsection*{You'll want to know: \mw{कारखानो}}
\textbf{\mw{कारखानो}} — \emph{kārkhāno} — \textquotedblleft{}a factory\textquotedblright{}.

\begin{quote}
\mw{कारखानो} \mw{सहर} \mw{सूं} \mw{दूर} \mw{है।}
\end{quote}

\begin{quote}
\emph{kārkhāno sahar sū̃ dūr hai.}
\end{quote}

\begin{quote}
The factory is far from the city.
\end{quote}

A masculine noun.

\begin{cousinweb}
\textbf{\mw{कारखानो}} (\emph{kārkhāno}) is usually traced to Persian \emph{kārkhāna}, \textquotedblleft{}workshop\textquotedblright{}.
\end{cousinweb}

\begin{grammarlens}[title={What you've built}]
One more word for this chapter.
\end{grammarlens}

\subsection*{Your turn}
\begin{itemize}
\raggedright
  \item \emph{Say it:} \emph{kārkhāno}
  \item \emph{Say it:} \emph{kārkhāno}, once more
  \item \emph{Say it:} \emph{kārkhāno sahar sū̃ dūr hai.} — The factory is far from the city.
  \item \emph{Recall:} say the Marwadi for a hospital, then the Marwadi for a sub-district office, then say \emph{kārkhāno} again
\end{itemize}

\begin{tcolorbox}[breakable,colback=teal!4,colframe=teal!35!black,title={Before you move on}]
What does \mw{कारखानो} mean? (\textquotedblleft{}A factory\textquotedblright{}.) Say it once more.
\end{tcolorbox}
\section[\texorpdfstring{\mw{अदालत} (adālat)}{अदालत (adālat)}]{\mw{अदालत} (adālat) — a court of law}

\label{lesson:MW-C208-adalat}



\begin{quote}
Before the new one: say the Marwadi for a sub-district office, then the Marwadi for a factory.
\end{quote}

\subsection*{You'll want to know: \mw{अदालत}}
\textbf{\mw{अदालत}} — \emph{adālat} — \textquotedblleft{}a court of law\textquotedblright{}.

\begin{quote}
\mw{अदालत} \mw{तहसील} \mw{कनै} \mw{है।}
\end{quote}

\begin{quote}
\emph{adālat tahsīl kanai hai.}
\end{quote}

\begin{quote}
The court is near the tehsil office.
\end{quote}

A feminine noun.

\begin{cousinweb}
\textbf{\mw{अदालत}} (\emph{adālat}) is usually traced, through Persian, to Arabic \emph{ʿadāla}, \textquotedblleft{}justice\textquotedblright{}.
\end{cousinweb}

\begin{grammarlens}[title={What you've built}]
One more word for this chapter.
\end{grammarlens}

\subsection*{Your turn}
\begin{itemize}
\raggedright
  \item \emph{Say it:} \emph{adālat}
  \item \emph{Say it:} \emph{adālat}, once more
  \item \emph{Say it:} \emph{adālat tahsīl kanai hai.} — The court is near the tehsil office.
  \item \emph{Recall:} say the Marwadi for a sub-district office, then the Marwadi for a factory, then say \emph{adālat} again
\end{itemize}

\begin{tcolorbox}[breakable,colback=teal!4,colframe=teal!35!black,title={Before you move on}]
What does \mw{अदालत} mean? (\textquotedblleft{}A court of law\textquotedblright{}.) Say it once more.
\end{tcolorbox}
\section[\texorpdfstring{\mw{खुलणो} (khulṇo)}{खुलणो (khulṇo)}]{\mw{खुलणो} (khulṇo) — to come open; to open (of a shop or office)}

\label{lesson:MW-C208-khulno}



\begin{quote}
Before the new one: say the Marwadi for a factory, then the Marwadi for a court of law.
\end{quote}

\subsection*{You'll want to know: \mw{खुलणो}}
\textbf{\mw{खुलणो}} — \emph{khulṇo} — \textquotedblleft{}to come open; to open (of a shop or office)\textquotedblright{}.

\begin{quote}
\mw{दफ्तर} \mw{दिनूगै} \mw{खुलै} \mw{है।}
\end{quote}

\begin{quote}
\emph{daphtar dinūgai khulai hai.}
\end{quote}

\begin{quote}
The office opens early in the morning.
\end{quote}

Past (it …): \mw{वो} \mw{दफ्तर} \mw{खुल्यो}, \emph{vo daphtar khulyo}.

\begin{cousinweb}
\textbf{\mw{खुलणो}} (\emph{khulṇo}) is inherited from a reconstructed \emph{*khull-}, \textquotedblleft{}to be open\textquotedblright{}; CDIAL lists Old Marwari \emph{khulo}, \textquotedblleft{}open\textquotedblright{}, and Hindi \emph{khulnā}, \textquotedblleft{}to be open\textquotedblright{} (CDIAL 3945). \textbf{\mw{खोलणो}}, \textquotedblleft{}to open\textquotedblright{}, is its transitive partner.
\end{cousinweb}

\begin{grammarlens}[title={What you've built}]
One more word for this chapter.
\end{grammarlens}

\subsection*{Your turn}
\begin{itemize}
\raggedright
  \item \emph{Say it:} \emph{khulṇo}
  \item \emph{Say it:} \emph{khulṇo}, once more
  \item \emph{Say it:} \emph{daphtar dinūgai khulai hai.} — The office opens early in the morning.
  \item \emph{Recall:} say the Marwadi for a factory, then the Marwadi for a court of law, then say \emph{khulṇo} again
\end{itemize}

\begin{tcolorbox}[breakable,colback=teal!4,colframe=teal!35!black,title={Before you move on}]
What does \mw{खुलणो} mean? (\textquotedblleft{}To come open; to open (of a shop or office)\textquotedblright{}.) Say it once more.
\end{tcolorbox}
